\begin{proposition}[KT-KOHNEN-THETA-NONCONGRUENCE]
The normalized Kohnen plus projection of $\theta^{p+4}$ is not,
in general for primes $p\equiv3\pmod4$, a scalar multiple modulo $p$ of the Cohen Eisenstein
series $H_{(p+3)/2}$.  In particular, this fails at $p=31$.
Consequently the Frobenius identity
$\theta^{p+4}\equiv\theta(q^p)\theta^4\pmod p$ alone does
not supply the missing simultaneous unit in the automatic Jacobi
family, even after the genuine plus projection.
\end{proposition}

\begin{proof}
Write $\theta(\tau)=\sum_{a\in\mathbf Z}q^{a^2}$,
$\theta_o(\tau)=\sum_{a\ \mathrm{odd}}q^{a^2}$ and
$\theta^m=\sum_{n\ge0}r_m(n)q^n$.
For an odd integer $b\ge1$ put $m=2b+1$ and
$s=(-1)^{(b+1)/2}$.  We use the normalization
\[
 \mathcal P f=\frac12f+
 \frac{s}{2\sqrt2}\sum_{\nu=0}^3 f|_{m/2}\xi\widetilde A_\nu,
 \quad
 \xi=\left(\begin{pmatrix}4&1\\0&4\end{pmatrix},e^{\pi i/4}\right),
 \quad A_\nu=\begin{pmatrix}1&0\\4\nu&1\end{pmatrix}.
                                                        \tag{KTHT.1}
\]
Here $(f|_k(A,\phi))(\tau)=\phi(\tau)^{-2k}f(A\tau)$,
and $\widetilde A_\nu$ uses the usual theta automorphy factor.
This is the normalized plus operator in Matsusaka's thesis,
printed pp.~27--28.  It is $3/2$ times the idempotent
normalization; this nonzero rational scale has no effect on the
question of proportionality modulo $31$.

The four transformed theta functions in (KTHT.1), before taking
their $m$th powers, are
\[
 e^{-\pi i/4}\theta(\tau+1/4),\quad
 \sqrt2\theta_o(\tau),\quad
 e^{\pi i/4}\theta(\tau-1/4),\quad
 \sqrt2\theta(4\tau),                                  \tag{KTHT.2}
\]
in the order $\nu=0,1,2,3$.
Here are details controlling the phases.  For $\nu=2$, factor
\[
 \begin{pmatrix}3&1/4\\8&1\end{pmatrix}
 =\begin{pmatrix}3&1\\8&3\end{pmatrix}
  \begin{pmatrix}1&-1/4\\0&1\end{pmatrix};
\]
the theta multiplier of the integral matrix is
$(8/3)\epsilon_3^{-1}=i$.
For $A=\left(\begin{smallmatrix}1&0\\2&1\end{smallmatrix}\right)$,
Poisson summation gives
\[
 (2\tau+1)^{-1/2}\theta(A\tau)
     =\sum_{a\ \mathrm{odd}}q^{a^2/4}.
\]
Factoring the matrix for $\nu=1$ as
$A\left(\begin{smallmatrix}2&1/4\\0&1/2\end{smallmatrix}\right)$
and using $a^2\equiv1\pmod8$ gives its entry in (KTHT.2).
For $\nu=3$, factor it as
\[
 \begin{pmatrix}4&1/4\\12&1\end{pmatrix}
 =\begin{pmatrix}3&1\\8&3\end{pmatrix}
  \begin{pmatrix}0&-1/4\\4&1\end{pmatrix}
\]
and use $\theta(-1/(4z))=\sqrt{-2iz}\,\theta(z)$.
The square roots have their holomorphic branches fixed at
$\tau=iy$, $y>0$, so these identities extend to the upper
half-plane without a sign ambiguity.

For $n\equiv3\pmod4$, the sum of the first and third
terms of (KTHT.2), raised to the $m$th power, has coefficient
$s\sqrt2\,r_m(n)$.  The last term is supported at multiples
of $4$, and the second starts at $q^m$.  Therefore
\[
 [q^n]\mathcal P\theta^m=r_m(n)
 \quad(0<n<m,\ n\equiv3\pmod4).                        \tag{KTHT.3}
\]
For $p=31$ and $m=35$, Frobenius gives
$r_{35}(3)\equiv r_4(3)=32\equiv1$ and
$r_{35}(7)\equiv r_4(7)=64\equiv2\pmod{31}$.
On the other hand, the defining Bernoulli-polynomial sums give
\[
 H_{17}(3)=-\frac{B_{17,\chi_{-3}}}{17}\equiv26,
 \qquad
 H_{17}(7)=-\frac{B_{17,\chi_{-7}}}{17}\equiv16
 \pmod{31}.
\]
Hence the two coefficient vectors have determinant
$1\cdot16-2\cdot26\equiv26\ne0\pmod{31}$.
They cannot be scalar multiples.  The standalone certificate
\texttt{kohnen\_theta\_discriminator.py} computes both exact
Bernoulli values and independently checks them using conductor
$31d$ character sums modulo $31^2$, for $d=3,7$; it also
computes the theta coefficients by exact integer convolution.
This disproves the proposed uniform scalar identity, not the
existence of another modular or genus-average construction.
\end{proof}

\paragraph{Source and normalization audit.}
T.~Matsusaka, \emph{Fourier coefficients of polyharmonic weak
Maass forms}, Ph.D.~thesis, Kyushu University (2019), printed
pp.~27--28, gives (KTHT.1), its slash convention, and the
phase $e^{\pi i/4}$; Remark~2.3.2 explicitly warns against
inserting an additional factor $2\lambda+1$ into that exponent
under this slash convention.  The institutional PDF is
\url{https://catalog.lib.kyushu-u.ac.jp/opac_download_md/2236044/math0246.pdf}.
The operator's plus-space property is an imported statement;
the theta transformations and the coefficient calculation above
are proved here.  No whole-thesis review is claimed.
The author PDF of Bruinier--Jenkins--Ono,
\emph{Hilbert class polynomials and traces of singular moduli},
Math.~Ann.~334 (2006), 373--393, \S3.2, formula~(3.12),
printed pp.~14--15, uses the idempotent factors $1/3$ instead of
$1/2$ but prints that extra phase factor.  We do not silently
use that inconsistent phase with its stated slash convention.
